% =============================================================
%  Section 1 — Introduction & Contexte
% =============================================================
\section{Contexte \& Vision stratégique}

\subsection{Positionnement}

\textbf{Jlax Expertise} est un cabinet indépendant à deux associés, issu du monde de l'infogérance franchisée, dont la mission est d'accompagner les PME dans leur modernisation technologique de manière \textbf{invisible, automatique et intégrée}.

\begin{pointcle}
  Notre différenciateur : la modernisation se déroule \emph{en arrière-plan}. Le client utilise ses outils habituels (Outlook, Teams, Odoo, Excel) et constate des gains sans rien avoir à changer.
\end{pointcle}

\subsection{Cible client}

\begin{itemize}
  \item PME de \textbf{5 à 50 employés}, tous secteurs
  \item Infrastructure existante : Microsoft 365, Odoo ou CRM équivalent
  \item Profil décideur : dirigeant non-technique, sensible au ROI et à la sécurité
  \item Douleurs exprimées : perte de temps, tâches répétitives, crainte cyber, image technologique vieillissante
\end{itemize}

\subsection{Périmètre d'intervention}

Quatre piliers couverts par notre offre :

\begin{enumerate}
  \item \textbf{Automatisation intelligente} — workflows invisibles connectant les outils existants
  \item \textbf{Intégration IA opérationnelle} — assistants IA métier, génération automatique de documents
  \item \textbf{Cybersécurité pragmatique} — SIEM managé, simulation phishing, score de sécurité
  \item \textbf{Pilotage \& reporting} — dashboards dirigeant consolidant KPIs métier et sécurité
\end{enumerate}

\subsection{Principe de modernisation invisible}

\begin{table}[htbp]
  \centering
  \caption{Ce que voit le client vs ce que nous déployons}
  \renewcommand{\arraystretch}{1.4}
  \begin{tabularx}{\textwidth}{>{\bfseries}l X X}
    \toprule
    \rowcolor{tabheader}
    \color{white}Domaine & \color{white}Ce que voit le client & \color{white}Ce que nous déployons \\
    \midrule
    Emails & Ses emails triés, résumés, sans action & IA classification via Graph API + n8n \\
    \rowcolor{rowgray}
    Sécurité & Un rapport mensuel dans sa boîte & Wazuh + Secure Score API + PDF auto \\
    CRM & Fiches complètes automatiquement & GPT-4o + hooks Odoo + n8n \\
    \rowcolor{rowgray}
    Devis & Devis pré-rempli depuis un email & Graph API → n8n → Odoo \\
    Reporting & Email hebdo avec ses KPIs & Metabase + Grafana + n8n scheduler \\
    \rowcolor{rowgray}
    Questions métier & Une réponse contextuelle Teams & Flowise RAG sur ses documents \\
    \bottomrule
  \end{tabularx}
\end{table}

\subsection{Contraintes structurantes}

\begin{attention}
  Tout outil nécessitant une formation, un changement d'interface, ou l'apprentissage d'un nouveau logiciel est \textbf{exclu du périmètre client}. Il peut exister en arrière-plan mais ne doit jamais être visible.
\end{attention}

Les trois règles non négociables :
\begin{itemize}
  \item \textbf{Zéro nouvel outil} à ouvrir pour le client
  \item \textbf{Zéro formation lourde} — adoption immédiate ou transparente
  \item \textbf{ROI mesurable} dès les 3 premiers mois (temps gagné, incidents évités, score amélioré)
\end{itemize}
